% GENERATED by diagnostics/build_rc_tables.R - do not hand-edit.
\begin{table}[htbp]\centering
\begin{threeparttable}
\caption{Where the break falls, by specification}
\begin{tabular}{lcccl}
\toprule
Specification & Break & sup-$F$ & Perm.\ $p$ & Bai--Perron \\
\midrule
\textbf{Count, verdict-bearing extracts} (chosen) & 2015 & 16.13$^{***}$ & 0.0002 & 1 break at 2015 \\
Count, all extracts (42-year basis) & 2015 & 12.33$^{**}$ & 0.0010 & 1 break at 2015 \\
Original v2 extract (Manchester spine) & 2015 & 19.05$^{***}$ & 0.0002 & 1 break at 2015 \\
Proportion, verdict-bearing extracts & 2015 & 6.14 & 0.0726 & 1 break at 2015 \\
Proportion, all extracts & 2017 & 7.65$^{*}$ & 0.0274 & 1 break at 2017 \\
Failure within the body's own control & 2015 & 4.02 & 0.2970 & none selected \\
\bottomrule
\end{tabular}
\begin{tablenotes}[flushleft]\footnotesize
\item The question this table answers is where the break falls, not whether it is significant.
\item 5 of the 6 specifications place the break in 2015, including the v2 extract used for the
conference presentation that preceded this work --- a different instrument, applied to a corpus
since expanded by about a tenth.
\item 2 of those 5 cannot distinguish the break from noise at conventional levels; they are shown
with their $p$ values rather than omitted. The agreement is in the date, not in the strength of
the evidence for it.
\item The chosen specification is the count of serious findings estimated from verdict-bearing
extracts. The alternatives are plotted in Figure~\ref{fig:ch02naoalt}.
\item Leave-one-out on the chosen series returns the same break year in 40 of 41 runs.
\item $^{*}p<.05$; $^{**}p<.01$; $^{***}p<.001$, permutation test.
\item Source: \texttt{diagnostics/fig214\_supporting\_tests.R}.
\end{tablenotes}
\end{threeparttable}
\label{tab:ch02naobreaks}
\end{table}
